\section{Strong anharmonicity: SrTiO\textsubscript{3}}
\label{sec:res_strong}
% ======================================================================

The static corrections of \cref{sec:res_scp} are a small perturbation for the
weakly anharmonic silicon systems. To test the opposite, strongly anharmonic
regime, where the harmonic approximation fails, we computed
FC2/FC3/FC4 for cubic SrTiO$_3$, the canonical
self-consistent-phonon (SCP) test case from~\cite{tadano_first-principles_2018}.

Cubic SrTiO$_3$ has unstable harmonic phonons: a ferroelectric (FE) mode at $\Gamma$ and an
antiferrodistortive (AFD) octahedral-rotation mode at
$R=(\tfrac12,\tfrac12,\tfrac12)$~\cite{tadano_first-principles_2018}. Fitting FC2 to
small-amplitude ($0.025$~\AA) rattled displacements in the same $2\times2\times2$
($40$-atom) cell used by Tadano \& Tsuneyuki recovers the AFD instability:
% REVIEW(verified): srtio3_rattle_renorm.py on the on-disk 60-structure fit gives omega_R = -1.656 THz = 55.2i cm^-1.
% REVIEW(open): the earlier "55i--92i, converging upward with the number of training structures" range has no surviving on-disk structure-count sweep (hiphive_fit.json holds a single 60-structure fit); dropped until the sweep is reproduced.
$\omega_R = 55i$~cm$^{-1}$ (the $60$-structure fit).
Their PBEsol value is $76i$~cm$^{-1}$. The static loop$+$tadpole magnitude is
% REVIEW(open): the 39 THz^2 value has no surviving on-disk calculation (REPRODUCE.md M23, pending); the ~6x is quoted against the top of the bulk-Si range.
$\|\Sigma_\mathrm{static}\|\approx39$~THz$^2$ at $300$~K --- about $6\times$ the
upper end of the bulk-Si range ($1$--$7$~THz$^2$ across $100$--$600$~K,
\cref{sec:res_scp}).

Fitting the same cell to a larger ($0.08$~\AA) rattle amplitude stabilises the AFD mode
% REVIEW(verified): the on-disk large-rattle fit (fc3_hiphive_srtio3_vasp) gives omega_R = +1.164 THz = +38.8 cm^-1; the previously quoted +90 cm^-1 belongs to a fit whose data no longer exists. Corrected to the reproducible value.
($\omega_R = +39$~cm$^{-1}$, \cref{fig:res_srtio3}). This is not a fitting error,
but the central identity of the effective-potential methods: a finite-amplitude fit
returns the effective, renormalised FC2 at the temperature where
$\langle u^2\rangle^{1/2}$ matches the rattle amplitude. Fitting at the thermal
amplitude is in the same spirit as the temperature-dependent effective potential
(TDEP)~\cite{hellman_temperature_2013} and the stochastic self-consistent harmonic
approximation (SSCHA)~\cite{monacelli_stochastic_2021}, though a single
fixed-amplitude static rattle is distinct from the thermal-ensemble sampling of
TDEP and the variational quantum self-consistent Gaussian of SSCHA. The difference between
the bare and effective FC2 is the static loop self-energy
$\tfrac12\Phi_4\!:\!\langle uu\rangle$ (\cref{sec:res_scp}).

\begin{figure}[htb]
  \centering
  \includegraphics[width=0.7\linewidth]{transport_sweeps/srtio3_rattle_renorm.pdf}
  \caption{SrTiO$_3$ phonon dispersion: FC2 fitted to a small ($0.025$~\AA,
  orange) versus large ($0.08$~\AA, blue) rattle amplitude.}
  \label{fig:res_srtio3}
\end{figure}

The renormalised phonons can hence be obtained either from the bare FC2 plus an
explicit quartic loop (the deterministic SCP of \cref{sec:res_scp}, requiring
FC4), or directly by fitting the effective FC2 at the thermal amplitude
(TDEP/SSCHA, no FC4 needed)~\cite{hellman_temperature_2013,monacelli_stochastic_2021}.
The explicit loop is needed when an analytic temperature dependence from a single
FC4 is wanted. The effective-fit route is preferred when FC4 is too costly.
For the weakly anharmonic silicon systems neither matters (effective $\approx$ bare FC2).
For SrTiO$_3$ one must use one route or the other.

A tensor rank-sweep of the device FC3 (truncated-SVD / HOSVD / CP) shows SrTiO$_3$ is
markedly more compressible than \texttt{d5a}. A truncated SVD reaches a
$7.6\%$ Frobenius error at rank $16$ for SrTiO$_3$ versus $18\%$ for \texttt{d5a}.

\paragraph{Anharmonic transport at $\eta=0$.}
Despite the strong anharmonicity, the SrTiO$_3$ transport device (built on a
$3\times3\times3$-supercell fit whose dispersion is stable --- no imaginary
modes, lowest optical branch at $4.4$~THz as $k_z\to0$) converges at $\eta=0$
like the carbon nanotube, with the default recipe and no soft-mode
stabilisation (\cref{fig:res_srtio3_transport}). At two cells the
three-phonon reduction is $G_\mathrm{anh}/G_\mathrm{ball}=0.728$ at $300$~K
(lead balance $1.2\times10^{-5}$, $45$ iterations) and deepens with length to
$0.627$ at three cells ($2.0\times10^{-4}$, $46$ iterations); across
temperature the ratio falls from $0.777$ at $200$~K to $0.728$ at $300$~K. The
$600$~K point did not close the convergence gate within the iteration budget
(best ratio $0.676$, lead balance $6.6\times10^{-3}$) and is reported as
unconverged --- consistent with the three-phonon phase space of the
perovskite growing rapidly with temperature. At matched length and $\eta=0$
the $300$~K reduction ($27\%$ at two cells) sits between the (3,3) tube's
($43\%$) and the (8,0) tube's ($22\%$).

\begin{figure}[htb]
  \centering
  \includegraphics[width=0.95\linewidth]{transport_sweeps/srtio3_eta0_transport.pdf}
  \caption{SrTiO$_3$ anharmonic transport at $\eta=0$ (production solver,
  $3\times3\times3$-supercell fit). \emph{Left}: ballistic and anharmonic conductance versus
  temperature at two cells; the hollow marker ($600$~K) did not converge and
  shows the best iterate. \emph{Right}: the ratio deepens with device length
  at $300$~K.}
  \label{fig:res_srtio3_transport}
\end{figure}
